% ======================================================================
% Chapter 2: The Klein-Gordon Field
% (book pages 13-34; PDF pages 36-57)
% ======================================================================
\chapter{The Klein-Gordon Field}

\section{The Necessity of the Field Viewpoint}

Quantum field theory is the application of quantum mechanics to dynamical
systems of fields, in the same sense that the basic course in quantum
mechanics is concerned mainly with the quantization of dynamical systems
of particles.  It is a subject that is absolutely essential for
understanding the current state of elementary particle physics.  With some
modification, the methods we will discuss also play a crucial role in the
most active areas of atomic, nuclear, and condensed-matter physics.  In
Part~I of this book, however, our primary concern will be with elementary
particles, and hence relativistic fields.

Given that we wish to understand processes that occur at very small
(quantum-mechanical) scales and very large (relativistic) energies, one
might still ask why we must study the quantization of fields.  Why can't
we just quantize relativistic particles the way we quantized
nonrelativistic particles?  This question can be answered on a number of
levels.  Perhaps the best approach is to write down a single-particle
relativistic wave equation (such as the Klein--Gordon equation or the
Dirac equation) and see that it gives rise to negative-energy states and
other inconsistencies.  Since this discussion usually takes place near the
end of a graduate-level quantum mechanics course, we will not repeat it
here.  It is easy, however, to understand why such an approach cannot
work.  We have no right to assume that any relativistic process can be
explained in terms of a single particle, since the Einstein relation
$E=mc^2$ allows for the creation of particle--antiparticle pairs.  Even
when there is not enough energy for pair creation, multiparticle states
appear, for example, as intermediate states in second-order perturbation
theory.  We can think of such states as existing only for a very short
time, according to the uncertainty principle $\Delta E\cdot\Delta t =
\hbar$.  As we go to higher orders in perturbation theory, arbitrarily many
such ``virtual'' particles can be created.

The necessity of having a multiparticle theory also arises in a less
obvious way, from considerations of causality.  Consider the amplitude for
a free particle to propagate from $\mathbf{x}_0$ to $\mathbf{x}$:
\begin{equation}
  U(t) = \langle \mathbf{x}|\ e^{-iHt}\ |\mathbf{x}_0\rangle .
\end{equation}
In nonrelativistic quantum mechanics we have $E = \mathbf{p}^2/2m$, so
\begin{equation}
  U(t) = \langle \mathbf{x}|\ e^{-i(\mathbf{p}^2/2m)t}\ |\mathbf{x}_0\rangle
  = \int \frac{d^3p}{(2\pi)^3}\ \langle \mathbf{x}|
    e^{-i(\mathbf{p}^2/2m)t}|\mathbf{p}\rangle\langle \mathbf{p}|\mathbf{x}_0\rangle
  = \frac{1}{(2\pi)^3} \int d^3p\ e^{-i(\mathbf{p}^2/2m)t}
    e^{i\mathbf{p}\cdot(\mathbf{x}-\mathbf{x}_0)}
  = \Bigl( \frac{m}{2\pi it} \Bigr)^{3/2}
    e^{im(\mathbf{x}-\mathbf{x}_0)^2/2t}.
\end{equation}
This expression is nonzero for all $x$ and $t$, indicating that a particle
can propagate between any two points in an arbitrarily short time.  In a
relativistic theory, this conclusion would signal a violation of
causality.  One might hope that using the relativistic expression
$E=\sqrt{\mathbf{p}^2+m^2}$ would help, but it does not.  In analogy with
the nonrelativistic case, we have
\begin{align}
  U(t) &= \langle \mathbf{x}|\ e^{-i\sqrt{\mathbf{p}^2+m^2}\,t}
    |\mathbf{x}_0\rangle
  = \frac{1}{(2\pi)^3} \int d^3p\
    e^{-i\sqrt{\mathbf{p}^2+m^2}\,t}\ e^{i\mathbf{p}\cdot(\mathbf{x}-\mathbf{x}_0)}\\
  &= \frac{1}{2\pi^2|\mathbf{x}-\mathbf{x}_0|}
    \int_0^\infty dp\ p\sin\bigl(p|\mathbf{x}-\mathbf{x}_0|\bigr)\
    e^{-i\sqrt{p^2+m^2}\,t}.
\end{align}
This integral can be evaluated explicitly in terms of Bessel functions.%
\footnote{See Gradshteyn and Ryzhik (1980), \#3.914.}  We will content
ourselves with looking at its asymptotic behavior for $x^2 > t^2$ (well
outside the light-cone), using the method of stationary phase.  The phase
function $px - t\sqrt{p^2+m^2}$ has a stationary point at
$p = imx/\sqrt{x^2-t^2}$.  We may freely push the contour upward so that
it goes through this point.  Plugging in this value for $p$, we find that,
up to a rational function of $x$ and $t$,
\begin{equation}
  U(t) \sim e^{-m\sqrt{x^2-t^2}} .
\end{equation}
Thus the propagation amplitude is small but nonzero outside the
light-cone, and causality is still violated.

Quantum field theory solves the causality problem in a miraculous way,
which we will discuss in Section~2.4.  We will find that, in the
multiparticle field theory, the propagation of a particle across a
spacelike interval is indistinguishable from the propagation of an
antiparticle in the opposite direction (see Fig.~\ref{fig:2.1}).  When we
ask whether an observation made at point $\mathbf{x}_0$ can affect an
observation made at point $\mathbf{x}$, we will find that the amplitudes
for particle and antiparticle propagation exactly cancel---so causality is
preserved.

\begin{figure}[t]
  \centering
  \includegraphics[width=0.5\textwidth]{fig2.1.png}
  \caption{Propagation from $x_0$ to $x$ in one frame looks like
    propagation from $x$ to $x_0$ in another frame.}
  \label{fig:2.1}
\end{figure}

Quantum field theory provides a natural way to handle not only
multiparticle states, but also transitions between states of different
particle number.  It solves the causality problem by introducing
antiparticles, then goes on to explain the relation between spin and
statistics.  But most important, it provides the tools necessary to
calculate innumerable scattering cross sections, particle lifetimes, and
other observable quantities.  The experimental confirmation of these
predictions, often to an unprecedented level of accuracy, is our real
reason for studying quantum field theory.

\section{Elements of Classical Field Theory}

In this section we review some of the formalism of classical field theory
that will be necessary in our subsequent discussion of quantum field
theory.

\subsection*{Lagrangian Field Theory}

The fundamental quantity of classical mechanics is the action, $S$, the
time integral of the Lagrangian, $L$.  In a local field theory the
Lagrangian can be written as the spatial integral of a Lagrangian density,
denoted by $\mathcal{L}$, which is a function of one or more fields
$\phi(x)$ and their derivatives $\partial_\mu\phi$.  Thus we have
\begin{equation}
  S = \int L\ dt = \int \mathcal{L}(\phi, \partial_\mu\phi)\ d^4x .
  \label{eq:2.1}
\end{equation}
Since this is a book on field theory, we will refer to $\mathcal{L}$
simply as the Lagrangian.

The principle of least action states that when a system evolves from one
given configuration to another between times $t_1$ and $t_2$, it does so
along the ``path'' in configuration space for which $S$ is an extremum
(normally a minimum).  We can write this condition as
\begin{equation}
  0 = \delta S = \int d^4x\ \Bigl\{
    \frac{\partial\mathcal{L}}{\partial\phi}\delta\phi
    + \frac{\partial\mathcal{L}}{\partial(\partial_\mu\phi)}
      \partial_\mu(\delta\phi) \Bigr\}.
  \label{eq:2.2}
\end{equation}
The last term can be turned into a surface integral over the boundary of
the four-dimensional spacetime region of integration.  Since the initial
and final field configurations are assumed given, $\delta\phi$ is zero at
the temporal beginning and end of this region.  If we restrict our
consideration to deformations $\delta\phi$ that vanish on the spatial
boundary of the region as well, then the surface term is zero.  Factoring
out the $\delta\phi$ from the first two terms, we note that, since the
integral must vanish for arbitrary $\delta\phi$, the quantity that
multiplies $\delta\phi$ must vanish at all points.  Thus we arrive at the
Euler--Lagrange equation of motion for a field,
\begin{equation}
  \partial_\mu \frac{\partial\mathcal{L}}{\partial(\partial_\mu\phi)}
  - \frac{\partial\mathcal{L}}{\partial\phi} = 0 .
  \label{eq:2.3}
\end{equation}
If the Lagrangian contains more than one field, there is one such equation
for each.

\subsection*{Hamiltonian Field Theory}

The Lagrangian formulation of field theory is particularly suited to
relativistic dynamics because all expressions are explicitly Lorentz
invariant.  Nevertheless we will use the Hamiltonian formulation
throughout the first part of this book, since it will make the transition
to quantum mechanics easier.  Recall that for a discrete system one can
define a conjugate momentum $p = \partial L/\partial\dot{q}$ (where
$\dot{q} = dq/dt$) for each dynamical variable $q$.  The Hamiltonian is
then $H = \sum p\dot{q} - L$.  The generalization to a continuous system
is best understood by pretending that the spatial points $\mathbf{x}$ are
discretely spaced.  We can define
\begin{equation}
  \pi(\mathbf{x}) \equiv \frac{\partial\mathcal{L}}
    {\partial\dot{\phi}(\mathbf{x})}
  = \frac{\delta}{\delta\dot{\phi}(\mathbf{x})}
    \int d^3y\ \mathcal{L}(\phi(y), \dot{\phi}(y))
  \approx \frac{\partial}{\partial\dot{\phi}(\mathbf{x})}
    \sum_{\mathbf{y}} \mathcal{L}(\phi(\mathbf{y}), \dot{\phi}(\mathbf{y}))
    \ d^3y = \frac{\partial\mathcal{L}}{\partial\dot{\phi}(\mathbf{x})}\ d^3x,
  \label{eq:2.4}
\end{equation}
where
\begin{equation}
  \pi(x) = \frac{\partial\mathcal{L}}{\partial\dot{\phi}(x)}
\end{equation}
is called the \emph{momentum density} conjugate to $\phi(x)$.  Thus the
Hamiltonian can be written
\begin{equation}
  H = \sum_{\mathbf{x}} \pi(\mathbf{x})\dot{\phi}(\mathbf{x}) - L .
\end{equation}
Passing to the continuum, this becomes
\begin{equation}
  H = \int d^3x\ \bigl[ \pi(x)\dot{\phi}(x) - \mathcal{L} \bigr]
    = \int d^3x\ \mathcal{H} .
  \label{eq:2.5}
\end{equation}
We will rederive this expression for the Hamiltonian density $\mathcal{H}$
near the end of this section, using a different method.

As a simple example, consider the theory of a single field $\phi(x)$,
governed by the Lagrangian
\begin{equation}
  \mathcal{L} = \tfrac{1}{2}\dot{\phi}^2 - \tfrac{1}{2}(\nabla\phi)^2
    - \tfrac{1}{2}m^2\phi^2 .
  \label{eq:2.6}
\end{equation}
For now we take $\phi$ to be a real-valued field.  The quantity $m$ will
be interpreted as a mass in Section~2.3, but for now just think of it as a
parameter.  From this Lagrangian the usual procedure gives the equation of
motion
\begin{equation}
  \Bigl( \frac{\partial^2}{\partial t^2} - \nabla^2 + m^2 \Bigr)\phi = 0
  \qquad\text{or}\qquad (\partial_\mu\partial^\mu + m^2)\phi = 0,
  \label{eq:2.7}
\end{equation}
which is the well-known Klein--Gordon equation.  (In this context it is a
classical field equation, like Maxwell's equations---not a
quantum-mechanical wave equation.)  Noting that the canonical momentum
density conjugate to $\phi(x)$ is $\pi(x) = \dot{\phi}(x)$, we can also
construct the Hamiltonian:
\begin{equation}
  H = \int d^3x\ \mathcal{H}
    = \int d^3x\ \bigl[ \tfrac{1}{2}\pi^2 + \tfrac{1}{2}(\nabla\phi)^2
      + \tfrac{1}{2}m^2\phi^2 \bigr].
  \label{eq:2.8}
\end{equation}
We can think of the three terms, respectively, as the energy cost of
``moving'' in time, the energy cost of ``shearing'' in space, and the
energy cost of having the field around at all.  We will investigate this
Hamiltonian much further in Sections~2.3 and~2.4.

\subsection*{Noether's Theorem}

Next let us discuss the relationship between symmetries and conservation
laws in classical field theory, summarized in Noether's theorem.  This
theorem concerns continuous transformations on the fields $\phi$, which in
infinitesimal form can be written
\begin{equation}
  \phi(x) \to \phi'(x) = \phi(x) + \alpha\Delta\phi(x),
  \label{eq:2.9}
\end{equation}
where $\alpha$ is an infinitesimal parameter and $\Delta\phi$ is some
deformation of the field configuration.  We call this transformation a
\emph{symmetry} if it leaves the equations of motion invariant.  This is
insured if the action is invariant under \eqref{eq:2.9}.  More generally,
we can allow the action to change by a surface term, since the presence of
such a term would not affect our derivation of the Euler--Lagrange
equations of motion \eqref{eq:2.3}.  The Lagrangian, therefore, must be
invariant under \eqref{eq:2.9} up to a 4-divergence:
\begin{equation}
  \mathcal{L}(x) \to \mathcal{L}(x) + \alpha\partial_\mu\mathcal{J}^\mu(x),
  \label{eq:2.10}
\end{equation}
for some $\mathcal{J}^\mu$.  Let us compare this expectation for
$\Delta\mathcal{L}$ to the result obtained by varying the fields:
\begin{equation}
  \Delta\mathcal{L}
  = \frac{\partial\mathcal{L}}{\partial\phi}\Delta\phi
    + \frac{\partial\mathcal{L}}{\partial(\partial_\mu\phi)}
      \partial_\mu(\Delta\phi)
  = \partial_\mu\Bigl(
      \frac{\partial\mathcal{L}}{\partial(\partial_\mu\phi)}
      \Delta\phi \Bigr).
  \label{eq:2.11}
\end{equation}
The second term vanishes by the Euler--Lagrange equation \eqref{eq:2.3}.
We set the remaining term equal to $\alpha\partial_\mu\mathcal{J}^\mu$ and
find
\begin{equation}
  \partial_\mu j^\mu(x) = 0,
  \qquad\text{for}\qquad
  j^\mu(x) = \frac{\partial\mathcal{L}}{\partial(\partial_\mu\phi)}
    \Delta\phi - \mathcal{J}^\mu .
  \label{eq:2.12}
\end{equation}
(If the symmetry involves more than one field, the first term of this
expression for $j^\mu(x)$ should be replaced by a sum of such terms, one
for each field.)  This result states that the current $j^\mu(x)$ is
conserved.  For each continuous symmetry of $\mathcal{L}$, we have such a
conservation law.

The conservation law can also be expressed by saying that the charge
\begin{equation}
  Q = \int_{\text{all space}} j^0\ d^3x
  \label{eq:2.13}
\end{equation}
is a constant in time.  Note, however, that the formulation of field
theory in terms of a local Lagrangian density leads directly to the local
form of the conservation law, Eq.~\eqref{eq:2.12}.

The easiest example of such a conservation law arises from a Lagrangian
with only a kinetic term: $\mathcal{L} =
\tfrac{1}{2}(\partial_\mu\phi)^2$.  The transformation $\phi\to\phi+\alpha$,
where $\alpha$ is a constant, leaves $\mathcal{L}$ unchanged, so we
conclude that the current $j^\mu = \partial^\mu\phi$ is conserved.  As a
less trivial example, consider the Lagrangian
\begin{equation}
  \mathcal{L} = |\partial_\mu\phi|^2 - m^2|\phi|^2,
  \label{eq:2.14}
\end{equation}
where $\phi$ is now a complex-valued field.  You can easily show that the
equation of motion for this Lagrangian is again the Klein--Gordon
equation, \eqref{eq:2.7}.  This Lagrangian is invariant under the
transformation $\phi \to e^{i\alpha}\phi$; for an infinitesimal
transformation we have
\begin{equation}
  \alpha\Delta\phi = i\alpha\phi;\qquad
  \alpha\Delta\phi^{*} = -i\alpha\phi^{*}.
  \label{eq:2.15}
\end{equation}
(We treat $\phi$ and $\phi^{*}$ as independent fields.  Alternatively, we
could work with the real and imaginary parts of $\phi$.)  It is now a
simple matter to show that the conserved Noether current is
\begin{equation}
  j^\mu = i\bigl[ (\partial^\mu\phi^{*})\phi - \phi^{*}(\partial^\mu\phi)
    \bigr].
  \label{eq:2.16}
\end{equation}
(The overall constant has been chosen arbitrarily.)  You can check
directly that the divergence of this current vanishes by using the
Klein--Gordon equation.  Later we will add terms to this Lagrangian that
couple $\phi$ to an electromagnetic field.  We will then interpret $j^\mu$
as the electromagnetic current density carried by the field, and the
spatial integral of $j^0$ as its electric charge.

Noether's theorem can also be applied to spacetime transformations such as
translations and rotations.  We can describe the infinitesimal translation
$x^\mu \to x^\mu - a^\mu$ alternatively as a transformation of the field
configuration
\begin{equation}
  \phi(x) \to \phi(x+a) = \phi(x) + a^\mu\partial_\mu\phi(x).
\end{equation}
The Lagrangian is also a scalar, so it must transform in the same way:
\begin{equation}
  \mathcal{L} \to \mathcal{L} + a^\mu\partial_\mu\mathcal{L}
    = \mathcal{L} + a^\nu\partial_\mu(\delta^\mu_\nu\mathcal{L}).
\end{equation}
Comparing this equation to \eqref{eq:2.10}, we see that we now have a
nonzero $\mathcal{J}^\mu$.  Taking this into account, we can apply the
theorem to obtain four separately conserved currents:
\begin{equation}
  T^{\mu\nu} = \frac{\partial\mathcal{L}}{\partial(\partial_\mu\phi)}
    \partial^\nu\phi - g^{\mu\nu}\mathcal{L} .
  \label{eq:2.17}
\end{equation}
This is precisely the stress-energy tensor, also called the
energy-momentum tensor, of the field $\phi$.  The conserved charge
associated with time translations is the Hamiltonian:
\begin{equation}
  H = \int T^{00}\ d^3x = \int \mathcal{H}\ d^3x.
  \label{eq:2.18}
\end{equation}
By computing this quantity for the Klein--Gordon field, one can recover
the result \eqref{eq:2.8}.  The conserved charges associated with spatial
translations are
\begin{equation}
  P^i = \int T^{0i}\ d^3x = \int \pi\partial^i\phi\ d^3x ,
  \label{eq:2.19}
\end{equation}
and we naturally interpret this as the (physical) momentum carried by the
field (not to be confused with the canonical momentum).

\section{The Klein-Gordon Field as Harmonic Oscillators}

We begin our discussion of quantum field theory with a rather formal
treatment of the simplest type of field: the real Klein--Gordon field.
The idea is to start with a classical field theory (the theory of a
classical scalar field governed by the Lagrangian \eqref{eq:2.6}) and then
``quantize'' it, that is, reinterpret the dynamical variables as operators
that obey canonical commutation relations.%
\footnote{This procedure is sometimes called \emph{second quantization}, to
distinguish the resulting Klein--Gordon equation (in which $\phi$ is an
operator) from the old one-particle Klein--Gordon equation (in which
$\phi$ was a wavefunction).  In this book we never adopt the latter point
of view; we start with a classical equation (in which $\phi$ is a
classical field) and quantize it exactly once.}  We will then ``solve''
the theory by finding the eigenvalues and eigenstates of the Hamiltonian,
using the harmonic oscillator as an analogy.

The classical theory of the real Klein--Gordon field was discussed briefly
(but sufficiently) in the previous section; the relevant expressions are
given in Eqs.~\eqref{eq:2.6}, \eqref{eq:2.7}, and \eqref{eq:2.8}.  To
quantize the theory, we follow the same procedure as for any other
dynamical system: We promote $\phi$ and $\pi$ to operators, and impose
suitable commutation relations.  Recall that for a discrete system of one
or more particles the commutation relations are
\begin{equation}
  [q_i, p_j] = i\delta_{ij};\qquad
  [q_i, q_j] = [p_i, p_j] = 0.
\end{equation}
For a continuous system the generalization is quite natural; since
$\pi(\mathbf{x})$ is the momentum density, we get a Dirac delta function
instead of a Kronecker delta:
\begin{equation}
  [\phi(\mathbf{x}), \pi(\mathbf{y})] = i\delta^{(3)}(\mathbf{x}-\mathbf{y});
  \qquad
  [\phi(\mathbf{x}), \phi(\mathbf{y})]
  = [\pi(\mathbf{x}), \pi(\mathbf{y})] = 0 .
  \label{eq:2.20}
\end{equation}
(For now we work in the Schr\"odinger picture where $\phi$ and $\pi$ do
not depend on time.  When we switch to the Heisenberg picture in the next
section, these ``equal time'' commutation relations will still hold
provided that both operators are considered at the same time.)

The Hamiltonian, being a function of $\phi$ and $\pi$, also becomes an
operator.  Our next task is to find the spectrum from the Hamiltonian.
Since there is no obvious way to do this, let us seek guidance by writing
the Klein--Gordon equation in Fourier space.  If we expand the classical
Klein--Gordon field as
\begin{equation}
  \phi(\mathbf{x}, t) = \int \frac{d^3p}{(2\pi)^3}\ e^{i\mathbf{p}\cdot\mathbf{x}}
    \phi(\mathbf{p}, t)
\end{equation}
(with $\phi^{*}(\mathbf{p}) = \phi(-\mathbf{p})$ so that $\phi(\mathbf{x})$
is real), the Klein--Gordon equation \eqref{eq:2.7} becomes
\begin{equation}
  \Bigl( \frac{\partial^2}{\partial t^2}
    + (|\mathbf{p}|^2 + m^2) \Bigr)\phi(\mathbf{p}, t) = 0 .
  \label{eq:2.21}
\end{equation}
This is the same as the equation of motion for a simple harmonic oscillator
with frequency
\begin{equation}
  \omega_{\mathbf{p}} = \sqrt{|\mathbf{p}|^2 + m^2}.
  \label{eq:2.22}
\end{equation}
The simple harmonic oscillator is a system whose spectrum we already know
how to find.  Let us briefly recall how it is done.  We write the
Hamiltonian as
\begin{equation}
  H_{\mathrm{SHO}} = \tfrac{1}{2}p^2 + \tfrac{1}{2}\omega^2\phi^2 .
\end{equation}
To find the eigenvalues of $H_{\mathrm{SHO}}$, we write $\phi$ and $p$ in
terms of ladder operators:
\begin{equation}
  \phi = \frac{1}{\sqrt{2\omega}}(a + a^\dagger);\qquad
  p = -i\sqrt{\frac{\omega}{2}}(a - a^\dagger) .
  \label{eq:2.23}
\end{equation}
The canonical commutation relation $[\phi, p] = i$ is equivalent to
\begin{equation}
  [a, a^\dagger] = 1.
  \label{eq:2.24}
\end{equation}
The Hamiltonian can now be rewritten
\begin{equation}
  H_{\mathrm{SHO}} = \omega\bigl( a^\dagger a + \tfrac{1}{2} \bigr).
\end{equation}
The state $|0\rangle$ such that $a|0\rangle = 0$ is an eigenstate of $H$
with eigenvalue $\tfrac{1}{2}\omega$, the zero-point energy.  Furthermore,
the commutators
\begin{equation}
  [H_{\mathrm{SHO}}, a^\dagger] = \omega a^\dagger,
  \qquad [H_{\mathrm{SHO}}, a] = -\omega a
\end{equation}
make it easy to verify that the states
\begin{equation}
  |n\rangle = (a^\dagger)^n |0\rangle
\end{equation}
are eigenstates of $H_{\mathrm{SHO}}$ with eigenvalues $(n+\tfrac{1}{2})
\omega$.  These states exhaust the spectrum.

We can find the spectrum of the Klein--Gordon Hamiltonian using the same
trick, but now each Fourier mode of the field is treated as an independent
oscillator with its own $a$ and $a^\dagger$.  In analogy with
\eqref{eq:2.23} we write
\begin{align}
  \phi(\mathbf{x}) &= \int \frac{d^3p}{(2\pi)^3}
    \frac{1}{\sqrt{2E_{\mathbf{p}}}}
    \bigl( a_{\mathbf{p}}e^{i\mathbf{p}\cdot\mathbf{x}}
      + a_{\mathbf{p}}^\dagger e^{-i\mathbf{p}\cdot\mathbf{x}} \bigr);
  \label{eq:2.25}\\
  \pi(\mathbf{x}) &= \int \frac{d^3p}{(2\pi)^3}
    (-i)\sqrt{\frac{E_{\mathbf{p}}}{2}}
    \bigl( a_{\mathbf{p}}e^{i\mathbf{p}\cdot\mathbf{x}}
      - a_{\mathbf{p}}^\dagger e^{-i\mathbf{p}\cdot\mathbf{x}} \bigr).
  \label{eq:2.26}
\end{align}
The inverse expressions for $a_{\mathbf{p}}$ and $a_{\mathbf{p}}^\dagger$
in terms of $\phi$ and $\pi$ are easy to derive but rarely needed.  In the
calculations below we will find it useful to rearrange \eqref{eq:2.25} and
\eqref{eq:2.26} as follows:
\begin{align}
  a_{\mathbf{p}} &= \int d^3x\ e^{-i\mathbf{p}\cdot\mathbf{x}}
    \Bigl[ \sqrt{\frac{E_{\mathbf{p}}}{2}}\phi(\mathbf{x})
      + \frac{i}{\sqrt{2E_{\mathbf{p}}}}\pi(\mathbf{x}) \Bigr];
  \label{eq:2.27}\\
  a_{\mathbf{p}}^\dagger &= \int d^3x\ e^{i\mathbf{p}\cdot\mathbf{x}}
    \Bigl[ \sqrt{\frac{E_{\mathbf{p}}}{2}}\phi(\mathbf{x})
      - \frac{i}{\sqrt{2E_{\mathbf{p}}}}\pi(\mathbf{x}) \Bigr].
  \label{eq:2.28}
\end{align}
The commutation relation \eqref{eq:2.24} becomes
\begin{equation}
  [a_{\mathbf{p}}, a_{\mathbf{p}'}^\dagger]
    = (2\pi)^3\delta^{(3)}(\mathbf{p}-\mathbf{p}'),
  \label{eq:2.29}
\end{equation}
from which you can verify that the commutator of $\phi$ and $\pi$ works
out correctly:
\begin{equation}
  [\phi(\mathbf{x}), \pi(\mathbf{x}')]
  = \int \frac{d^3p\ d^3p'}{(2\pi)^6}\ \frac{-i}{2}
    e^{i(\mathbf{p}\cdot\mathbf{x}+\mathbf{p}'\cdot\mathbf{x}')}
  = i\delta^{(3)}(\mathbf{x}-\mathbf{x}').
  \label{eq:2.30}
\end{equation}
(If computations such as this one and the next are unfamiliar to you,
please work them out carefully; they are quite easy after a little
practice, and are fundamental to the formalism of the next two chapters.)

We are now ready to express the Hamiltonian in terms of ladder operators.
Starting from its expression \eqref{eq:2.8} in terms of $\phi$ and $\pi$,
we have
\begin{equation}
  H = \int \frac{d^3p}{(2\pi)^3}\ E_{\mathbf{p}}
    \Bigl( a_{\mathbf{p}}^\dagger a_{\mathbf{p}}
      + \tfrac{1}{2}[a_{\mathbf{p}}, a_{\mathbf{p}}^\dagger] \Bigr)
  = \int \frac{d^3p}{(2\pi)^3}\ E_{\mathbf{p}}
    \Bigl( a_{\mathbf{p}}^\dagger a_{\mathbf{p}}
      + \tfrac{1}{2}(2\pi)^3\delta^{(3)}(0) \Bigr).
  \label{eq:2.31}
\end{equation}
The second term is proportional to $\delta^{(3)}(0)$, an infinite c-number.
It is simply the sum over all modes of the zero-point energies
$\omega_{\mathbf{p}}/2$, so its presence is completely expected, if
somewhat disturbing.  Fortunately, this infinite energy shift cannot be
detected experimentally, since experiments measure only energy differences
from the ground state of $H$.  We will therefore ignore this infinite
constant term in all of our calculations.  It is possible that this energy
shift of the ground state could create a problem at a deeper level in the
theory; we will discuss this matter in the Epilogue.

Using this expression for the Hamiltonian in terms of $a_{\mathbf{p}}$ and
$a_{\mathbf{p}}^\dagger$, it is easy to evaluate the commutators
\begin{equation}
  [H, a_{\mathbf{p}}^\dagger] = \omega_{\mathbf{p}}a_{\mathbf{p}}^\dagger,
  \qquad [H, a_{\mathbf{p}}] = -\omega_{\mathbf{p}}a_{\mathbf{p}} .
  \label{eq:2.32}
\end{equation}
We can now write down the spectrum of the theory, just as for the harmonic
oscillator.  The state $|0\rangle$ such that $a_{\mathbf{p}}|0\rangle = 0$
for all $\mathbf{p}$ is the ground state or \emph{vacuum}, and has $E=0$
after we drop the infinite constant in \eqref{eq:2.31}.  All other energy
eigenstates can be built by acting on $|0\rangle$ with creation operators.
In general, the state
$a_{\mathbf{p}}^\dagger a_{\mathbf{q}}^\dagger\cdots|0\rangle$ is an
eigenstate of $H$ with energy $\omega_{\mathbf{p}} + \omega_{\mathbf{q}}
+ \cdots$.  These states exhaust the spectrum.

Having found the spectrum of the Hamiltonian, let us try to interpret its
eigenstates.  From \eqref{eq:2.19} and a calculation similar to
\eqref{eq:2.31} we can write down the total momentum operator,
\begin{equation}
  \mathbf{P} = -\int d^3x\ \pi(\mathbf{x})\nabla\phi(\mathbf{x})
  = \int \frac{d^3p}{(2\pi)^3}\ \mathbf{p}\ a_{\mathbf{p}}^\dagger
    a_{\mathbf{p}} .
  \label{eq:2.33}
\end{equation}
So the operator $a_{\mathbf{p}}^\dagger$ creates momentum $\mathbf{p}$ and
energy $\omega_{\mathbf{p}} = \sqrt{|\mathbf{p}|^2+m^2}$.  Similarly, the
state $a_{\mathbf{p}}^\dagger a_{\mathbf{q}}^\dagger\cdots|0\rangle$ has
momentum $\mathbf{p}+\mathbf{q}+\cdots$.  It is quite natural to call
these excitations \emph{particles}, since they are discrete entities that
have the proper relativistic energy-momentum relation.  (By a particle we
do not mean something that must be localized in space;
$a_{\mathbf{p}}^\dagger$ creates particles in momentum eigenstates.)  From
now on we will refer to $\omega_{\mathbf{p}}$ as $E_{\mathbf{p}}$ (or
simply $E$), since it really is the energy of a particle.  Note, by the
way, that the energy is always positive:
$E_{\mathbf{p}} = +\sqrt{|\mathbf{p}|^2+m^2}$.

This formalism also allows us to determine the statistics of our
particles.  Consider the two-particle state
$a_{\mathbf{p}}^\dagger a_{\mathbf{q}}^\dagger|0\rangle$.  Since
$a_{\mathbf{p}}^\dagger$ and $a_{\mathbf{q}}^\dagger$ commute, this state
is identical to the state $a_{\mathbf{q}}^\dagger a_{\mathbf{p}}^\dagger
|0\rangle$ in which the two particles are interchanged.  Moreover, a
single mode $\mathbf{p}$ can contain arbitrarily many particles (just as a
simple harmonic oscillator can be excited to arbitrarily high levels).
Thus we conclude that Klein--Gordon particles obey Bose--Einstein
statistics.

We naturally choose to normalize the vacuum state so that
$\langle 0|0\rangle = 1$.  The one-particle states
$|\mathbf{p}\rangle \propto a_{\mathbf{p}}^\dagger|0\rangle$ will also
appear quite often, and it is worthwhile to adopt a convention for their
normalization.  The simplest normalization
$\langle\mathbf{p}|\mathbf{q}\rangle = (2\pi)^3\delta^{(3)}
(\mathbf{p}-\mathbf{q})$ (which many books use) is not Lorentz invariant,
as we can demonstrate by considering the effect of a boost in the
3-direction.  Under such a boost we have
$p^{3'} = \gamma(p^3+\beta E)$, $E' = \gamma(E+\beta p^3)$.  Using the
delta function identity
\begin{equation}
  \delta\bigl( f(x) - f(x_0) \bigr) =
    \frac{1}{|f'(x_0)|}\,\delta(x-x_0),
  \label{eq:2.34}
\end{equation}
we can compute
\begin{align}
  \delta^{(3)}(\mathbf{p}-\mathbf{q})
  &= \delta^{(3)}(\mathbf{p}'-\mathbf{q}')\
    \Bigl|\frac{\partial p^3}{\partial p^{3'}}\Bigr|
  = \delta^{(3)}(\mathbf{p}'-\mathbf{q}')\
    \gamma\Bigl(1+\beta\frac{\partial E}{\partial p^3}\Bigr)\\
  &= \delta^{(3)}(\mathbf{p}'-\mathbf{q}')\
    \gamma\Bigl(1+\beta\frac{p^3}{E}\Bigr)
  = \delta^{(3)}(\mathbf{p}'-\mathbf{q}')\
    \frac{\gamma}{E}(E+\beta p^3)
  = \delta^{(3)}(\mathbf{p}'-\mathbf{q}')\ \frac{E'}{E} .
\end{align}
The problem is that volumes are not invariant under boosts; a box whose
volume is $V$ in its rest frame has volume $V/\gamma$ in a boosted frame,
due to Lorentz contraction.  But from the above calculation, we see that
the quantity $E_{\mathbf{p}}\,\delta^{(3)}(\mathbf{p}-\mathbf{q})$ is
Lorentz invariant.  We therefore define
\begin{equation}
  |\mathbf{p}\rangle = \sqrt{2E_{\mathbf{p}}}\ a_{\mathbf{p}}^\dagger |0\rangle,
  \label{eq:2.35}
\end{equation}
so that
\begin{equation}
  \langle\mathbf{p}|\mathbf{q}\rangle
    = 2E_{\mathbf{p}}\,(2\pi)^3\delta^{(3)}(\mathbf{p}-\mathbf{q}).
  \label{eq:2.36}
\end{equation}
(The factor of 2 is unnecessary, but is convenient because of the factor
of 2 in Eq.~\eqref{eq:2.25}.)

On the Hilbert space of quantum states, a Lorentz transformation $\Lambda$
will be implemented as some unitary operator $U(\Lambda)$.  Our
normalization condition \eqref{eq:2.35} then implies that
\begin{equation}
  U(\Lambda)\,|\mathbf{p}\rangle = |\Lambda\mathbf{p}\rangle .
  \label{eq:2.37}
\end{equation}
If we prefer to think of this transformation as acting on the operator
$a_{\mathbf{p}}^\dagger$, we can also write
\begin{equation}
  U(\Lambda)\,a_{\mathbf{p}}^\dagger\,U^{-1}(\Lambda)
    = a_{\Lambda\mathbf{p}}^\dagger .
  \label{eq:2.38}
\end{equation}
With this normalization we must divide by $2E_{\mathbf{p}}$ in other
places.  For example, the completeness relation for the one-particle
states is
\begin{equation}
  \int \frac{d^3p}{(2\pi)^3}\ \frac{1}{2E_{\mathbf{p}}}\
    |\mathbf{p}\rangle\langle\mathbf{p}| = 1 ,
  \label{eq:2.39}
\end{equation}
where the operator on the left is simply the identity within the subspace
of one-particle states, and zero in the rest of the Hilbert space.
Integrals of this form will occur quite often; in fact, the integral
\begin{equation}
  \int \frac{d^3p}{(2\pi)^3}\ \frac{1}{2E_{\mathbf{p}}}\ f(p)
  = \int \frac{d^4p}{(2\pi)^3}\ f(p)\ \delta(p^2-m^2)\Big|_{p^0>0}
  \label{eq:2.40}
\end{equation}
is a Lorentz-invariant 3-momentum integral, in the sense that if $f(p)$ is
Lorentz-invariant, so is $\int d^3p\ f(p)/(2E_{\mathbf{p}})$.  The
integration can be thought of as being over the $p^0>0$ branch of the
hyperboloid $p^2=m^2$ in 4-momentum space (see Fig.~\ref{fig:2.2}).

\begin{figure}[t]
  \centering
  \includegraphics[width=0.45\textwidth]{fig2.2.png}
  \caption{The Lorentz-invariant 3-momentum integral is over the upper
    branch of the hyperboloid $p^2=m^2$.}
  \label{fig:2.2}
\end{figure}

Finally let us consider the interpretation of the state $\phi(\mathbf{x})
|0\rangle$.  From the expansion \eqref{eq:2.25} we see that
\begin{equation}
  \phi(\mathbf{x})|0\rangle = \int \frac{d^3p}{(2\pi)^3}\
    \frac{1}{2E_{\mathbf{p}}}\ e^{-i\mathbf{p}\cdot\mathbf{x}}\
    |\mathbf{p}\rangle
  \label{eq:2.41}
\end{equation}
is a linear superposition of single-particle states that have well-defined
momentum.  Except for the factor $1/2E_{\mathbf{p}}$, this is the same as
the familiar nonrelativistic expression for the eigenstate of position
$|\mathbf{x}\rangle$; in fact the extra factor is nearly constant for small
(nonrelativistic) $p$.  We will therefore put forward the same
interpretation, and claim that the operator $\phi(\mathbf{x})$, acting on
the vacuum, creates a particle at position $\mathbf{x}$.  This
interpretation is further confirmed when we compute
\begin{equation}
  \langle 0|\phi(\mathbf{x})|\mathbf{p}\rangle
    = \langle 0| \int \frac{d^3p'}{(2\pi)^3}\
      \frac{1}{\sqrt{2E_{\mathbf{p}'}}}\ e^{-i\mathbf{p}'\cdot\mathbf{x}}\
      a_{\mathbf{p}'}\ \sqrt{2E_{\mathbf{p}}}\ a_{\mathbf{p}}^\dagger|0\rangle
    = e^{-i\mathbf{p}\cdot\mathbf{x}} .
  \label{eq:2.42}
\end{equation}
We can interpret this as the position-space representation of the
single-particle wavefunction of the state $|\mathbf{p}\rangle$, just as in
nonrelativistic quantum mechanics $\langle\mathbf{x}|\mathbf{p}\rangle
\propto e^{i\mathbf{p}\cdot\mathbf{x}}$ is the wavefunction of the state
$|\mathbf{p}\rangle$.

\section{The Klein-Gordon Field in Space-Time}

In the previous section we quantized the Klein--Gordon field in the
Schr\"odinger picture, and interpreted the resulting theory in terms of
relativistic particles.  In this section we will switch to the Heisenberg
picture, where it will be easier to discuss time-dependent quantities and
questions of causality.  After a few preliminaries, we will return to the
question of acausal propagation raised in Section~2.1.  We will also derive
an expression for the Klein--Gordon propagator, a crucial part of the
Feynman rules to be developed in Chapter~4.

In the Heisenberg picture, we make the operators $\phi$ and $\pi$
time-dependent in the usual way:
\begin{equation}
  \phi(\mathbf{x}) = \phi(x,t) = e^{iHt}\phi(\mathbf{x})e^{-iHt},
  \label{eq:2.43}
\end{equation}
and similarly for $\pi(\mathbf{x}) = \pi(x,t)$.  The Heisenberg equation
of motion,
\begin{equation}
  i\frac{\partial}{\partial t}\phi = [\phi, H],
  \label{eq:2.44}
\end{equation}
allows us to compute the time dependence of $\phi$ and $\pi$:
\begin{align}
  i\dot{\phi}(x,t) &= i\int d^3x'\ i\delta^{(3)}(\mathbf{x}-\mathbf{x}')\
    \pi(\mathbf{x}', t) = i\pi(x,t);\\
  i\dot{\pi}(x,t) &= i\int d^3x'\ \Bigl[ \tfrac{1}{2}(\nabla\phi)^2
    + \tfrac{1}{2}m^2\phi^2 \Bigr]
  = i(\nabla^2 - m^2)\phi(x,t).
\end{align}
Combining the two results gives
\begin{equation}
  \frac{\partial^2}{\partial t^2}\phi = (\nabla^2 - m^2)\phi,
  \label{eq:2.45}
\end{equation}
which is just the Klein--Gordon equation.

We can better understand the time dependence of $\phi(x)$ and $\pi(x)$ by
writing them in terms of creation and annihilation operators.  First note
that
\begin{equation}
  Ha_{\mathbf{p}} = a_{\mathbf{p}}(H - E_{\mathbf{p}}),
  \label{eq:2.46}
\end{equation}
and hence
\begin{equation}
  H^n a_{\mathbf{p}} = a_{\mathbf{p}}(H - E_{\mathbf{p}})^n,
\end{equation}
for any $n$.  A similar relation (with $-$ replaced by $+$) holds for
$a_{\mathbf{p}}^\dagger$.  Thus we have derived the identities
which we can use on expression \eqref{eq:2.25} for $\phi(\mathbf{x})$ to
find the desired expression for the Heisenberg operator $\phi(x)$,
according to \eqref{eq:2.43}.  (We will always use the symbols
$a_{\mathbf{p}}$ and $a_{\mathbf{p}}^\dagger$ to represent the
time-independent, Schr\"odinger-picture ladder operators.)  The result is
\begin{equation}
  \phi(x) = \int \frac{d^3p}{(2\pi)^3}\ \frac{1}{\sqrt{2E_{\mathbf{p}}}}
    \Bigl( a_{\mathbf{p}}e^{-ip\cdot x}
      + a_{\mathbf{p}}^\dagger e^{ip\cdot x} \Bigr).
  \label{eq:2.47}
\end{equation}
It is worth mentioning that we can perform the same manipulations with
$\mathbf{P}$ instead of $H$ to relate $\phi(x)$ to $\phi(0)$.  In analogy
with \eqref{eq:2.46}, one can show
\begin{equation}
  P^\mu a_{\mathbf{p}} = a_{\mathbf{p}}(P^\mu - p^\mu),
  \label{eq:2.48}
\end{equation}
and therefore
\begin{equation}
  \phi(x) = e^{i(P\cdot x - P^0 x^0)}\phi(0)e^{-i(P\cdot x - P^0x^0)}
  = e^{iP\cdot x}\phi(0)e^{-iP\cdot x},
  \label{eq:2.49}
\end{equation}
where $P^\mu = (H, \mathbf{P})$.  (The notation here is confusing but
standard.  Remember that $P^\mu$ is the momentum operator, whose eigenvalue
is the total momentum of the system.  On the other hand, $p^\mu$ is the
momentum of a single Fourier mode of the field, which we interpret as the
momentum of a particle in that mode.  For a one-particle state of
well-defined momentum, $p^\mu$ is the eigenvalue of $P^\mu$.)

Equation \eqref{eq:2.47} makes explicit the dual particle and wave
interpretations of the quantum field $\phi(x)$.  On the one hand,
$\phi(x)$ is written as a Hilbert space operator, which creates and
destroys the particles that are the quanta of field excitation.  On the
other hand, $\phi(x)$ is written as a linear combination of solutions
($e^{ip\cdot x}$ and $e^{-ip\cdot x}$) of the Klein--Gordon equation.
Both signs of the time dependence in the exponential appear: We find both
$e^{-ip\cdot x}$ and $e^{+ip\cdot x}$, although $p^0$ is always positive.
If these were single-particle wavefunctions, they would correspond to
states of positive and negative energy; let us refer to them more
generally as positive- and negative-frequency modes.  The connection
between the particle creation operators and the waveforms displayed here
is always valid for free quantum fields: A positive-frequency solution of
the field equation has as its coefficient the operator that destroys a
particle in that single-particle wavefunction.  A negative-frequency
solution of the field equation, being the Hermitian conjugate of a
positive-frequency solution, has as its coefficient the operator that
creates a particle in that positive-energy single-particle wavefunction.
In this way, the fact that relativistic wave equations have both
positive- and negative-frequency solutions is reconciled with the
requirement that a sensible quantum theory contain only positive
excitation energies.

\subsection{Causality}

Now let us return to the question of causality raised at the beginning of
this chapter.  In our present formalism, still working in the Heisenberg
picture, the amplitude for a particle to propagate from $y$ to $x$ is
$\langle 0|\phi(x)\phi(y)|0\rangle$.  We will call this quantity
$D(x-y)$.  Each operator $\phi$ is a sum of $a$ and $a^\dagger$ operators,
but only the term $\langle 0|a_{\mathbf{p}}a_{\mathbf{q}}^\dagger|0\rangle
= (2\pi)^3\delta^{(3)}(\mathbf{p}-\mathbf{q})$ survives in this
expression.  It is easy to check that we are left with
\begin{equation}
  D(x-y) = \langle 0|\phi(x)\phi(y)|0\rangle
  = \int \frac{d^3p}{(2\pi)^3}\ \frac{1}{2E_{\mathbf{p}}}\
    e^{-ip\cdot(x-y)} .
  \label{eq:2.50}
\end{equation}
We have already argued in \eqref{eq:2.40} that integrals of this form are
Lorentz invariant.  Let us now evaluate this integral for some particular
values of $x-y$.  First consider the case where the difference $x-y$ is
purely in the time-direction: $x^0-y^0=t$, $\mathbf{x}-\mathbf{y}=0$.
(If the interval from $y$ to $x$ is timelike, there is always a frame in
which this is the case.)  Then we have
\begin{align}
  D(x-y) &= \int_0^\infty \frac{dp}{(2\pi)^3}\ \frac{p^2}{2\sqrt{p^2+m^2}}\
    e^{-i\sqrt{p^2+m^2}\,t}
  = \frac{1}{4\pi^2}\int_m^\infty dE\ \sqrt{E^2-m^2}\ e^{-iEt}\\
  &\sim \frac{m}{\pi^2}\ e^{-imt},\qquad t\to\infty .
  \label{eq:2.51}
\end{align}
Next consider the case where $x-y$ is purely spatial: $x^0-y^0=0$,
$\mathbf{x}-\mathbf{y}=\mathbf{r}$.  The amplitude is then
\begin{align}
  D(x-y) &= \int \frac{d^3p}{(2\pi)^3}\ \frac{1}{2E_{\mathbf{p}}}\
    e^{i\mathbf{p}\cdot\mathbf{r}}
  = \frac{1}{(2\pi)^3}\int_0^\infty dp\ \frac{p^2}{2E_p}\
    \frac{e^{ipr}-e^{-ipr}}{ipr}\\
  &= \frac{1}{4\pi^2 i r}\int_{-\infty}^{\infty} dp\ \frac{p\ e^{ipr}}
    {2\sqrt{p^2+m^2}} .
\end{align}
The integrand, considered as a complex function of $p$, has branch cuts on
the imaginary axis starting at $\pm im$ (see Fig.~\ref{fig:2.3}).  To
evaluate the integral we push the contour up to wrap around the upper
branch cut.  Defining $\rho = -ip$, we obtain
\begin{equation}
  D(x-y) = \frac{1}{4\pi^2 r}\int_m^\infty d\rho\
    \frac{\rho\ e^{-\rho r}}{\sqrt{\rho^2-m^2}}
  \sim \frac{e^{-mr}}{4\pi r},\qquad r\to\infty .
  \label{eq:2.52}
\end{equation}
So again we find that outside the light-cone, the propagation amplitude is
exponentially vanishing but nonzero.

\begin{figure}[t]
  \centering
  \includegraphics[width=0.5\textwidth]{fig2.3.png}
  \caption{Contour for evaluating propagation amplitude $D(x-y)$ over a
    spacelike interval.}
  \label{fig:2.3}
\end{figure}

To really discuss causality, however, we should ask not whether particles
can propagate over spacelike intervals, but whether a measurement
performed at one point can affect a measurement at another point whose
separation from the first is spacelike.  The simplest thing we could try
to measure is the field $\phi(x)$, so we should compute the commutator
$[\phi(x), \phi(y)]$; if this commutator vanishes, one measurement cannot
affect the other.  In fact, if the commutator vanishes for
$(x-y)^2<0$, causality is preserved quite generally, since commutators
involving any function of $\phi(x)$, including $\pi(x) = \partial\phi/
\partial t$, would also have to vanish.  Of course we know from Eq.
\eqref{eq:2.20} that the commutator vanishes for $x^0=y^0$; now let's do
the more general computation:
\begin{equation}
  [\phi(x), \phi(y)] = D(x-y) - D(y-x).
  \label{eq:2.53}
\end{equation}
When $(x-y)^2<0$, we can perform a Lorentz transformation on the second
term (since each term is separately Lorentz invariant), taking
$(x-y)\to-(x-y)$, as shown in Fig.~\ref{fig:2.4}.  The two terms are
therefore equal and cancel to give zero; causality is preserved.  Note
that if $(x-y)^2>0$ there is no continuous Lorentz transformation that
takes $(x-y)\to-(x-y)$.  In this case, by Eq.~\eqref{eq:2.51}, the
amplitude is (fortunately) nonzero, roughly $(e^{-imt}-e^{+imt})$ for the
special case $\mathbf{x}-\mathbf{y}=0$.  Thus we conclude that no
measurement in the Klein--Gordon theory can affect another measurement
outside the light-cone.  Causality is maintained in the Klein--Gordon
theory just as suggested at the end of Section~2.1.  To understand this
mechanism properly, however, we should broaden the context of our
discussion to include a complex Klein--Gordon field, which has distinct
particle and antiparticle excitations.

\begin{figure}[t]
  \centering
  \includegraphics[width=0.5\textwidth]{fig2.4.png}
  \caption{When $x-y$ is spacelike, a continuous Lorentz transformation
    can take $(x-y)$ to $-(x-y)$.}
  \label{fig:2.4}
\end{figure}

As was mentioned in the discussion of Eq.~\eqref{eq:2.15}, we can add a
conserved charge to the Klein--Gordon theory by considering the field
$\phi(x)$ to be complex- rather than real-valued.  When the complex scalar
field theory is quantized (see Problem~2.2), $\phi(x)$ will create
positively charged particles and destroy negatively charged ones, while
$\phi^\dagger(x)$ will perform the opposite operations.  Then the
commutator $[\phi(x), \phi^\dagger(y)]$ will have nonzero contributions,
which must delicately cancel outside the light-cone to preserve causality.
The two contributions have the spacetime interpretation of the two terms
in \eqref{eq:2.53}, but with charges attached.  The first term will
represent the propagation of a negatively charged particle from $y$ to
$x$.  The second term will represent the propagation of a positively
charged particle from $x$ to $y$.  In order for these two processes to be
present and give canceling amplitudes, both of these particles must exist,
and they must have the same mass.  In quantum field theory, then, causality
requires that every particle have a corresponding antiparticle with the
same mass and opposite quantum numbers (in this case electric charge).  For
the real-valued Klein--Gordon field, the particle is its own antiparticle.

\subsection{The Klein-Gordon Propagator}

Let us study the commutator $[\phi(x), \phi(y)]$ a little further.  Since
it is a c-number, we can write $[\phi(x), \phi(y)] = \langle 0|[\phi(x),
\phi(y)]|0\rangle$.  This can be rewritten as a four-dimensional integral
as follows, assuming for now that $x^0>y^0$:
\begin{align}
  \langle 0|[\phi(x), \phi(y)]|0\rangle
  &= \int \frac{d^3p}{(2\pi)^3}\
    \frac{1}{2E_{\mathbf{p}}}\Bigl( e^{-ip\cdot(x-y)} - e^{ip\cdot(x-y)}
      \Bigr)\\
  &= \int \frac{d^3p}{(2\pi)^3}\ \frac{1}{2E_{\mathbf{p}}}\
    \Bigl( e^{-ip\cdot(x-y)}\Big|_{p^0=E_{\mathbf{p}}}
      - e^{-ip\cdot(x-y)}\Big|_{p^0=-E_{\mathbf{p}}} \Bigr)\\
  &= \int \frac{d^3p}{(2\pi)^3}\ \int \frac{dp^0}{2\pi i}\
    \frac{e^{-ip\cdot(x-y)}}{p^2-m^2}
  \label{eq:2.54}
\end{align}
In the last step the $p^0$ integral is to be performed along the following
contour:
\[
  \begin{array}{cc}
    -E_{\mathbf{p}} & +E_{\mathbf{p}}
  \end{array}
\]
For $x^0>y^0$ we can close the contour below, picking up both poles to
obtain the previous line of \eqref{eq:2.54}.  For $x^0<y^0$ we may close
the contour above, giving zero.  Thus the last line of \eqref{eq:2.54},
together with the prescription for going around the poles, is an
expression for what we will call
\begin{equation}
  D_R(x-y) = \theta(x^0-y^0)\ \langle 0|[\phi(x), \phi(y)]|0\rangle .
  \label{eq:2.55}
\end{equation}
To understand this quantity better, let's do another computation:
\begin{align}
  (\partial^2 + m^2)D_R(x-y)
  &= \bigl(\partial^2\theta(x^0-y^0)\bigr)\ \langle 0|[\phi(x), \phi(y)]|0\rangle\\
  &\quad + 2\bigl(\partial_\mu\theta(x^0-y^0)\bigr)\
    \bigl(\partial^\mu\langle 0|[\phi(x), \phi(y)]|0\rangle\bigr)\\
  &\quad + \theta(x^0-y^0)\ (\partial^2+m^2)\ \langle 0|[\phi(x), \phi(y)]|0\rangle\\
  &= -\dot{\delta}(x^0-y^0)\ \langle 0|[\phi(x), \phi(y)]|0\rangle
    + 2\delta(x^0-y^0)\ \langle 0|[\pi(x), \phi(y)]|0\rangle + 0\\
  &= -i\delta^{(4)}(x-y).
  \label{eq:2.56}
\end{align}
This says that $D_R(x-y)$ is a Green's function of the Klein--Gordon
operator.  Since it vanishes for $x^0<y^0$, it is the \emph{retarded}
Green's function.

If we had not already derived expression \eqref{eq:2.54}, we could find it
by Fourier transformation.  Writing
\begin{equation}
  D_R(x-y) = \int \frac{d^4p}{(2\pi)^4}\ e^{-ip\cdot(x-y)}\ D_R(p),
  \label{eq:2.57}
\end{equation}
we obtain an algebraic expression for $D_R(p)$:
\begin{equation}
  (-p^2+m^2)D_R(p) = -i .
\end{equation}
Thus we immediately arrive at the result
\begin{equation}
  D_R(p) = \frac{i}{p^2-m^2} .
  \label{eq:2.58}
\end{equation}
The $p^0$-integral of \eqref{eq:2.58} can be evaluated according to four
different contours, of which that used in \eqref{eq:2.54} is only one.  In
Chapter~4 we will find that a different pole prescription,
\[
  \begin{array}{cc}
    -E_{\mathbf{p}} & +E_{\mathbf{p}}
  \end{array}
\]
is extremely useful; it is called the \emph{Feynman prescription}.  A
convenient way to remember it is to write
\begin{equation}
  D_F(x-y) = \int \frac{d^4p}{(2\pi)^4}\ \frac{i}{p^2-m^2+i\epsilon}\
    e^{-ip\cdot(x-y)}
  \label{eq:2.59}
\end{equation}
since the poles are then at $p^0 = \pm(E_{\mathbf{p}}-i\epsilon)$,
displaced properly above and below the real axis.  When $x^0>y^0$ we can
perform the $p^0$ integral by closing the contour below, obtaining exactly
the propagation amplitude $D(x-y)$ \eqref{eq:2.50}.  When $x^0<y^0$ we
close the contour above, obtaining the same expression but with $x$ and
$y$ interchanged.  Thus we have
\begin{equation}
  D_F(x-y) = \begin{cases}
    D(x-y) & \text{for } x^0>y^0\\
    D(y-x) & \text{for } x^0<y^0
  \end{cases}
  = \theta(x^0-y^0)\langle 0|\phi(x)\phi(y)|0\rangle
    + \theta(y^0-x^0)\langle 0|\phi(y)\phi(x)|0\rangle
  = \langle 0|T\phi(x)\phi(y)|0\rangle .
  \label{eq:2.60}
\end{equation}
The last line defines the ``time-ordering'' symbol $T$, which instructs us
to place the operators that follow in order with the latest to the left.
By applying $(\partial^2+m^2)$ to the last line, you can verify directly
that $D_F$ is a Green's function of the Klein--Gordon operator.

Equations \eqref{eq:2.59} and \eqref{eq:2.60} are, from a practical point
of view, the most important results of this chapter.  The Green's function
$D_F(x-y)$ is called the \emph{Feynman propagator} for a Klein--Gordon
particle, since it is, after all, a propagation amplitude.  Indeed, the
Feynman propagator will turn out to be part of the Feynman rules:
$D_F(x-y)$ (or $D_F(p)$) is the expression that we will attach to internal
lines of Feynman diagrams, representing the propagation of virtual
particles.

Nevertheless we are still a long way from being able to do any real
calculations, since so far we have talked only about the free Klein--Gordon
theory, where the field equation is linear and there are no interactions.
Individual particles live in their isolated modes, oblivious to each
others' existence and to the existence of any other species of particles.
In such a theory there is no hope of making any observations, by scattering
or any other means.  On the other hand, the formalism we have developed is
extremely important, since the free theory forms the basis for doing
perturbative calculations in the interacting theory.

\subsection{Particle Creation by a Classical Source}

There is one type of interaction, however, that we are already equipped to
handle.  Consider a Klein--Gordon field coupled to an external, classical
source field $j(x)$.  That is, consider the field equation
\begin{equation}
  (\partial^2 + m^2)\phi(x) = j(x),
  \label{eq:2.61}
\end{equation}
where $j(x)$ is some fixed, known function of space and time that is
nonzero only for a finite time interval.  If we start in the vacuum state,
what will we find after $j(x)$ has been turned on and off again?

The field equation \eqref{eq:2.61} follows from the Lagrangian
\begin{equation}
  \mathcal{L} = \tfrac{1}{2}(\partial_\mu\phi)^2 - \tfrac{1}{2}m^2\phi^2
    + j(x)\phi(x) .
  \label{eq:2.62}
\end{equation}
But if $j(x)$ is turned on for only a finite time, it is easiest to solve
the problem using the field equation directly.  Before $j(x)$ is turned
on, $\phi(x)$ has the form
\[
  \phi_0(x) = \int \frac{d^3p}{(2\pi)^3}\ \frac{1}{\sqrt{2E_{\mathbf{p}}}}
    \Bigl( a_{\mathbf{p}}e^{-ip\cdot x}
      + a_{\mathbf{p}}^\dagger e^{ip\cdot x} \Bigr).
\]
If there were no source, this would be the solution for all time.  With a
source, the solution of the equation of motion can be constructed using
the retarded Green's function:
\begin{align}
  \phi(x) &= \phi_0(x) + i\int d^4y\ D_R(x-y)\,j(y)\\
  &= \phi_0(x) + \int \frac{d^3p}{(2\pi)^3}\ \frac{i}{\sqrt{2E_{\mathbf{p}}}}
    \Bigl( e^{-ip\cdot x}\tilde{j}(p) - e^{ip\cdot x}\tilde{j}^{*}(p)
      \Bigr),
  \label{eq:2.63}
\end{align}
where we have used $D_R(x-y)$ from Eq.~\eqref{eq:2.58}.  If we wait until
all of $j$ is in the past, the theta function equals 1 in the whole domain
of integration.  Then $\phi(x)$ involves only the Fourier transform of
$j$,
\begin{equation}
  \tilde{j}(p) = \int d^4y\ e^{ip\cdot y}\,j(y),
\end{equation}
evaluated at 4-momenta $p$ such that $p^2=m^2$.  It is natural to group
the positive-frequency terms together with $a_{\mathbf{p}}$ and the
negative-frequency terms with $a_{\mathbf{p}}^\dagger$; this yields the
expression
\begin{equation}
  \phi(x) = \int \frac{d^3p}{(2\pi)^3}\ \frac{1}{\sqrt{2E_{\mathbf{p}}}}
    \Bigl\{ \Bigl( a_{\mathbf{p}}
      + \frac{i\tilde{j}(p)}{\sqrt{2E_{\mathbf{p}}}} \Bigr)e^{-ip\cdot x}
      + \mathrm{h.c.} \Bigr\}.
  \label{eq:2.64}
\end{equation}
You can now guess (or compute) the form of the Hamiltonian after $j(x)$
has acted: Just replace $a_{\mathbf{p}}$ with
$(a_{\mathbf{p}} + i\tilde{j}(p)/\sqrt{2E_{\mathbf{p}}})$ to obtain
\begin{equation}
  H = \int \frac{d^3p}{(2\pi)^3}\ E_{\mathbf{p}}\
    \Bigl( a_{\mathbf{p}}^\dagger
      + \frac{-i\tilde{j}^{*}(p)}{\sqrt{2E_{\mathbf{p}}}} \Bigr)
    \Bigl( a_{\mathbf{p}}
      + \frac{i\tilde{j}(p)}{\sqrt{2E_{\mathbf{p}}}} \Bigr) .
\end{equation}
The energy of the system after the source has been turned off is
\begin{equation}
  \langle 0|H|0\rangle = \int \frac{d^3p}{(2\pi)^3}\
    \frac{|\tilde{j}(p)|^2}{2E_{\mathbf{p}}},
  \label{eq:2.65}
\end{equation}
where $|0\rangle$ still denotes the ground state of the free theory.  We
can interpret these results in terms of particles by identifying
$|\tilde{j}(p)|^2/2E_{\mathbf{p}}$ as the probability density for creating
a particle in the mode $\mathbf{p}$.  Then the total number of particles
produced is
\begin{equation}
  \langle 0|N|0\rangle = \int \frac{d^3p}{(2\pi)^3}\
    \frac{|\tilde{j}(p)|^2}{2E_{\mathbf{p}}},
  \label{eq:2.66}
\end{equation}
where $N = \int d^3p\ a_{\mathbf{p}}^\dagger a_{\mathbf{p}}/(2\pi)^3$.
Only those Fourier components of $j(x)$ that are in resonance with
on-mass-shell (i.e., $p^2=m^2$) Klein--Gordon waves are effective at
creating particles.  We will return to this subject in Problem~4.1.  In
Chapter~6 we will study the analogous problem of photon creation by an
accelerated electron (bremsstrahlung).

\problems

\textbf{2.1} \emph{Classical electromagnetism} (with no sources) follows
from the action
\begin{equation}
  S = -\frac{1}{4}\int d^4x\ F_{\mu\nu}F^{\mu\nu},
\end{equation}
where $F_{\mu\nu} = \partial_\mu A_\nu - \partial_\nu A_\mu$.
\begin{enumerate}[label=(\alph*)]
  \item Derive Maxwell's equations as the Euler--Lagrange equations of
    this action, treating the components $A^\mu(x)$ as the dynamical
    variables.  Write the equations in standard form by identifying
    $E^i = -F^{0i}$ and $\epsilon^{ijk}B^k = -F^{ij}$.
  \item Construct the energy-momentum tensor for this theory.  Note that
    the usual procedure does not result in a symmetric tensor.  To remedy
    that, we can add to $T^{\mu\nu}$ a term of the form
    $\partial_\lambda K^{\lambda\mu\nu}$, where $K^{\lambda\mu\nu}$ is
    antisymmetric in its first two indices.  Such an object is
    automatically divergenceless, so $T^{\mu\nu} = \cdots + \partial_\lambda
    K^{\lambda\mu\nu}$ is an equally good energy-momentum tensor with the
    same globally conserved energy and momentum.  Show that this
    construction, with
    \begin{equation}
      K^{\lambda\mu\nu} = F^{\mu\lambda}A^\nu,
    \end{equation}
    leads to an energy-momentum tensor $T^{\mu\nu}$ that is symmetric and
    yields the standard formulae for the electromagnetic energy and
    momentum densities:
    \begin{equation}
      \mathcal{E} = \tfrac{1}{2}(\mathbf{E}^2+\mathbf{B}^2);\qquad
      \mathbf{S} = \mathbf{E}\times\mathbf{B}.
    \end{equation}
\end{enumerate}

\textbf{2.2} \emph{The complex scalar field.}  Consider the field theory of
a complex-valued scalar field obeying the Klein--Gordon equation.  The
action of this theory is
\begin{equation}
  S = \int d^4x\ \bigl( \partial_\mu\phi^{*}\partial^\mu\phi - m^2\phi^{*}\phi
    \bigr).
\end{equation}
It is easiest to analyze this theory by considering $\phi(x)$ and
$\phi^{*}(x)$, rather than the real and imaginary parts of $\phi(x)$, as
the basic dynamical variables.
\begin{enumerate}[label=(\alph*)]
  \item Find the conjugate momenta to $\phi(x)$ and $\phi^{*}(x)$ and the
    canonical commutation relations.  Show that the Hamiltonian is
    \begin{equation}
      H = \int d^3x\ \bigl( \pi^{*}\pi + \nabla\phi^{*}\cdot\nabla\phi
        + m^2\phi^{*}\phi \bigr).
    \end{equation}
    Compute the Heisenberg equation of motion for $\phi(x)$ and show that
    it is indeed the Klein--Gordon equation.
  \item Diagonalize $H$ by introducing creation and annihilation
    operators.  Show that the theory contains two sets of particles of
    mass $m$.
  \item Rewrite the conserved charge
    \begin{equation}
      Q = \int d^3x\ \frac{i}{2}\bigl( \phi^{*}\pi^{*} - \pi\phi \bigr)
    \end{equation}
    in terms of creation and annihilation operators, and evaluate the
    charge of the particles of each type.
  \item Consider the case of two complex Klein--Gordon fields with the
    same mass.  Label the fields as $\phi_a(x)$, where $a=1,2$.  Show that
    there are now four conserved charges, one given by the generalization
    of part (c), and the other three given by
    \begin{equation}
      Q^l = \int d^3x\ \frac{i}{2}\bigl(
        \phi^{*}_a(\sigma^l)_{ab}\phi_b - \pi_a(\sigma^l)_{ab}\pi^{*}_b \bigr),
    \end{equation}
    where $\sigma^l$ are the Pauli sigma matrices.  Show that these three
    charges have the commutation relations of angular momentum (SU(2)).
    Generalize these results to the case of $n$ identical complex scalar
    fields.%
    \footnote{With some additional work you can show that there are
    actually six conserved charges in the case of two complex fields, and
    $n(2n-1)$ in the case of $n$ fields, corresponding to the generators
    of the rotation group in four and $2n$ dimensions, respectively.  The
    extra symmetries often do not survive when nonlinear interactions of
    the fields are included.}
\end{enumerate}

\textbf{2.3} Evaluate the function
\begin{equation}
  \langle 0|\phi(x)\phi(y)|0\rangle = D(x-y) = \int \frac{d^3p}{(2\pi)^3}
    \frac{1}{2E_{\mathbf{p}}}\ e^{-ip\cdot(x-y)},
\end{equation}
for $(x-y)$ spacelike so that $(x-y)^2 = -r^2$, explicitly in terms of
Bessel functions.
